\section{Setup and review of the Oaxaca--Blinder decomposition (OBD)}
\label{sec:counterfactual_OBD}

We first review the index number problem and OBD, defining them at the population level where we have full information about all variables of interest.
Then we describe how to proceed with real data.

%In this section, we revisit the OBD from a counterfactual perspective. 
%We show how this view naturally leads to index-number problem; i.e., that the choice of reference group in the OBD is arbitrary. 
%We first show how the index-number problem arises at the level of counterfactual mean comparisons, before introducing the standard linear formulation of the OBD.

\subsection{Differences in means as a counterfactual exercise} 
\label{sec:counterfactual_example}

As a first step toward describing the OBD, we illustrate the index number problem by framing it as a counterfactual exercise as in \citet{fortinLemieuxFirpo2011}.
% Unlike \citet{fortinLemieuxFirpo2011}, we observe that this counterfactual exercise need not invoke a linear model, which we leave off until \ref{sec:ob} below.
We consider an observed outcome $Y \in \mathbb{R}$ and covariates $X$ in some space $\mathcal{X}$. We take two groups, $H$ and $K$. We assume that the tuple $(X,Y)$ is independently and identically distributed within each group $g$ with corresponding expectation $\E_g$. We suppose we observe a gap in the mean outcomes, $\E_H[Y]$ and $\E_K[Y]$, and would like to understand to what extent the gap is due to differences in the distribution of $X$ versus differences in the distribution of $Y\mid X$ between the two groups. 
For the moment, we assume we have access to population-level quantities, including all expectations.

To analyze the gap between $\E_H[Y]$ and $\E_K[Y]$, we might construct a counterfactual driven only by differences in covariates.
Again taking our example, where $Y$ represents mortality and $X$ represents risk factors, we might ask:
(A) \emph{How would Hospital $H$'s mortality rate change if its patients had the risk factor distribution of patients at Hospital $K$?}
But, this question starts from the patient group at Hospital $H$ as a reference. We could instead start from group $K$ and ask:
(B) \emph{How would Hospital $K$'s mortality rate change if its patients had the risk factor distribution of patients at Hospital $H$?}

% Analogously, the gender wage gap forms a classic example \citep{oaxaca1973, blinder1973} with $Y$ as wages and $X$ as education. We might ask:
% (A) \emph{How much would women's ($H$) mean wages change if women kept their current remuneration relationship with education but had the educational distribution of men ($K$)?}
% or instead 
% (B) \emph{How much would men's ($K$) mean wages change if men kept their current remuneration relationship with education but had the educational distribution of women ($H$)?}

More generally, we might ask:
(A) \emph{How much would Group $H$'s mean outcome ($Y$) change if Group $H$ kept its outcome--covariate relationship ($Y \mid X$) but had the covariate ($X$) distribution of Group $K$?}
or 
(B) \emph{How much would Group $K$'s mean outcome ($Y$) change if Group $K$ kept its outcome--covariate relationship ($Y \mid X$) but had the covariate ($X$) distribution of Group $H$?}

To address the general questions, it will be useful to compute the mean of a counterfactual $Y$ that has the outcome--covariate relationship ($Y \mid X$) of group $g$ and the covariate ($X$) distribution of group $g'$. By the law of total expectation, this mean equals
$\E_{g'}[\E_g[Y \mid X]]$.
Then the difference in mean outcomes between the two groups can be decomposed into differences between each group's mean
and the chosen counterfactual mean. The decompositions corresponding to questions (A) and (B) above are, respectively, 
\cref{eq:OB_H_Counterfactual,eq:OB_K_Counterfactual} below.
%
% AISTATS two-column: reflowed so each decomposition sits on its own line (no words changed). Original kept below.
{\color{red}\small
\begin{align}
\Delta_Y &:= \E_H[Y] - \E_K[Y] \nonumber \\
    \label{eq:OB_H_Counterfactual}
    &=\underbrace{\E_H[Y] - \E_K[\E_H[Y \mid X]] }_{E^{(H)}} 
     + \underbrace{\E_K[\E_H[Y \mid X]] - \E_K[Y]}_{U^{(H)} } \\
     \label{eq:OB_K_Counterfactual}
     &=\underbrace{\E_H[Y] - \E_H[\E_K[Y \mid X]] }_{U^{(K)}} 
     + \underbrace{\E_H[\E_K[Y \mid X]] - \E_K[Y] }_{E^{(K)}}.
\end{align}}
%REV \begin{align}
%REV \Delta_Y := \E_H[Y] - \E_K[Y] 
%REV     \label{eq:OB_H_Counterfactual}
%REV     &=\underbrace{\E_H[Y] - \E_K[\E_H[Y \mid X]] }_{E^{(H)}} 
%REV      + \underbrace{\E_K[\E_H[Y \mid X]] - \E_K[Y]}_{U^{(H)} } \\
%REV      \label{eq:OB_K_Counterfactual}
%REV      &=\underbrace{\E_H[Y] - \E_H[\E_K[Y \mid X]] }_{U^{(K)}} 
%REV      + \underbrace{\E_H[\E_K[Y \mid X]] - \E_K[Y] }_{E^{(K)}}.
%REV \end{align} 
%Alternatively, the decomposition corresponding to question (B) above is:
%\begin{align}
%    \label{eq:OB_K_Counterfactual}
%    &\E_H[Y] - \E_K[Y]= \\
%    &\underbrace{\Big(\E_H[Y] - M_{H,K} \Big)}_{U^{(K)}} \nonumber 
%     + \underbrace{\Big(M_{H,K} - \E_K[Y]\Big)}_{E^{(K)}}.
%\end{align}
%
We follow the traditional naming and notation for these terms \citep{fortinLemieuxFirpo2011}.
The term corresponding to changing the covariates is called the
``explained component'' since this part of the difference is explained by the observed
covariates. The remaining term, corresponding to changing the conditional $Y \mid X$ distribution,
is called the ``unexplained component'' since it is not explained by the observed covariates. We denote these components as $E$ and $U$, respectively, with superscript denoting the reference group choice.

We might interpret $E$ and $U$ as the respective contributions of the covariates and outcome--covariate relationship to the 
difference in outcome means between groups $H$ and $K$. \textcolor{red}{This interpretation is descriptive; the OBD does not carry a causal interpretation without further assumptions. We refer readers to Section 2.3 of \citet{fortinLemieuxFirpo2011} for a formal treatment of causal identification for the OBD, and to \citet{quintas2024multiply,flachaire2025decomposing} for recent work.} The \textbf{index number problem}, then, represents the observation
that in general, $E^{(K)} \neq E^{(H)}$ and $U^{(K)} \neq U^{(H)}$.

\textbf{A conceptual resolution?}
Before proceeding, we observe that there is not an immediate and general conceptual resolution of the index number problem.
It is not clear to us in the hospital example above that one direction is more natural, or even that the two directions represent distinct policy questions.

Generic responses to the index number problem do not resolve it without relying on arbitrary choices or compromising the interpretation of results; see \cref{sec:existing_index_number} for detailed discussion, which we summarize here.
First, approaches like those of \citet{oaxacaRansom1998,neumark1988} define a new reference group by aggregating $H$ and $K$. However, there are many ways to perform this aggregation, introducing a new, still-arbitrary choice.
Second, use of Shapley values \citep{Shapley1953} could call for averaging the explained or unexplained components for each reference group.
However, we show that the choice of reference group can reverse the sign of either component, complicating their interpretation. 
Shapley-style approaches thus average two questionable quantities and obscure information when the averaged components have opposite signs.

In practice, we do not see a consensus approach, even within fields of study.
For instance, in economic analyses of wage gaps, some argue that a single reference group makes the most sense to use in context, as in \citet{dinardoFortinLemieux1996}'s study of unionization, the minimum wage, and economic inequality.
Often, the choice is a single reference group that aggregates the two original references in some way \citep{oaxacaRansom1994,reimers1983,cotton1988,kassenboehmer2014distributional,sharafRashad2016,rahimiNazari2021,Allen2022,bachan2022genderchancellor}.
However, we also observe published work that takes either men or women as the reference group \citep{nielsen2000wage,piazzalunga2019increase,leythienne2021paygapeu,tao2024gender,paradaContzen2025}.
Some papers report results for both references \citep{ONeill2006,fortinLemieuxFirpo2011,sharafRashad2016,rahimiNazari2021}. 
Still others report results for just one reference, with no argument for excluding the other reference \citep{zhang_2019, ayubiShahbaziKhazaei2024,paradaContzen2025,singleton2016,lamJamiesonMittinty2021,leythienne2021paygapeu,alHanawiNjagi2022,tao2024gender}.

% \input{Sections/02b_setup}

\subsection{Real-data Oaxaca--Blinder decomposition and drawing conclusions} 
\label{sec:ob_real}
In practice, researchers do not have access to population quantities and must instead draw conclusions from finite-data estimates.
In particular, implementing \cref{eq:OB_H_Counterfactual,eq:OB_K_Counterfactual} requires estimating $\E_g[Y]$, typically with the sample mean, and the counterfactual mean $\E_{g'}[\E_g[Y \mid X]]$, which requires more thought. 
Practitioners typically estimate counterfactual means by (1) using data from $g$ to compute an estimator of $\E_g[Y \mid X]$ (essentially a regression), and (2) using data from $g'$ to take the sample average across the estimator from the first step.
Concretely, for each $g \in {\{H, K\}}$, suppose we have $N_g$ observations of $\{X_n, Y_n\}$. Let $\hat{M}_g(X)$ be an estimator for $\E_g[Y \mid X]$.
Then the following equation implements \cref{eq:OB_H_Counterfactual}:
% AISTATS two-column: split over two lines; last term corrected to N_K (typo fixed in the workshop version). Original kept below.
{\color{red}
\begin{align}	
    \hat{\Delta}_Y
    \label{eq:NOBD_H_estimated}
    &=\underbrace{ \frac{1}{N_H} \sum_{n=1}^{N_H} Y_n - \frac{1}{N_K} \sum_{n=1}^{N_K} \hat{M}_H(X_n)}_{\hat{E}^{(H)}} \nonumber \\
    &\quad + \underbrace{\frac{1}{N_K} \sum_{n=1}^{N_K} \hat{M}_H(X_n) - \frac{1}{N_K} \sum_{n=1}^{N_K} Y_n}_{\hat{U}^{(H)} }.
\end{align}}
%REV \begin{align}	
%REV     \hat{\Delta}_Y
%REV     \label{eq:NOBD_H_estimated}
%REV     &=\underbrace{ \frac{1}{N_H} \sum_{n=1}^{N_H} Y_n - \frac{1}{N_K} \sum_{n=1}^{N_K} \hat{M}_H(X_n)}_{\hat{E}^{(H)}} 
%REV      + \underbrace{\frac{1}{N_K} \sum_{n=1}^{N_K} \hat{M}_H(X_n) - \frac{1}{N_H} \sum_{n=1}^{N_H} Y_n}_{\hat{U}^{(H)} }.
%REV \end{align}

In the literature, the term \emph{OBD} historically refers to using \cref{eq:OB_H_Counterfactual,eq:OB_K_Counterfactual} with the further choice of linear regression to compute $\hat{M}_g(X)$.
We emphasize that the use of linear regression is a modeling choice, and that it is possible to choose nonlinear estimators for $\E_g  [ Y | X ]$.
For instance, \citet{quintas2024multiply,flachaire2025decomposing,mao2025double} explore generalizations of the OBD that use machine learning methods to estimate $\E_g  [ Y | X ]$, and derive asymptotic properties for such procedures.
To distinguish from the standard OBD, we will refer to the case where $\hat{M}_g(X)$ is a nonlinear estimator as the \emph{nonlinear Oaxaca-Blinder decomposition} (NOBD).

\textbf{Drawing conclusions from the (N)OBD.}
It is common to draw substantive conclusions based on the signs (and statistical significance) of the empirical explained and unexplained components \citep{oaxacaRansom1998, jann2008,fortinLemieuxFirpo2011}. Returning to the hospital example from above, suppose Hospital $H$ has lower mortality ($\Delta_Y < 0$), and serves as the reference group. 
If the (N)OBD explained component has a negative sign, we might conclude that $H$ can attribute
its lower mortality rate in part to patient characteristics, such as healthier blood pressure. 
Analogously, if the (N)OBD unexplained component has a significant negative sign, we might attribute group $H$'s lower
mortality in part to the quality of hospital care.
In either case, we might also require that the observed components be statistically distinguishable from zero to form a conclusion; without significance, we might hesitate to attribute differences in mortality to covariates or care quality. 
For the rest of the paper, we will consider a practitioner drawing their conclusions based on the (statistically significant) sign of the explained or unexplained component.
